\section{Discussion}\label{sec:discussion}
\subsection{What constitutes a useful audit result}
Cryptographic validity establishes the specified integer relation. Observed
agreement describes a finite workload. An assignment certificate lower-bounds
correct labels, while a cost-change certificate can establish an update
decision without certifying every label. These are different outputs: 100\%
of requested update signs being resolved does not mean 100\% assignment
coverage. A valid inconclusive proof should not be presented as a successful
decision or as an exact-service speedup.

The sign-stratified grid contains 24 naturally worsening updates, and neither
representation certifies any of them directly. Requiring a negative upper
endpoint before approval defers these tasks; exact fallback is needed to learn
their true sign. This is selective improvement certification, not demonstrated
natural-regression detection. Early-training image improvements provide a
useful regime, but do not by themselves validate a broad deployed-model
monitoring claim.

The paired extension targets the decision gap left by overlapping absolute
intervals. Early image updates give positive examples, and the public fallback
demonstrates how inconclusive tasks can be resolved. However, mature updates
and text expose a substantial negative result: uncertainty often exceeds the
cost change. The original 1,024-record workload's 5.08-fold component saving
and 98.25\% agreement remain an approximation trade-off rather than evidence
of useful update ordering. New workloads and a changed decision interface do
not erase that limitation.

For fixed centers and repeated subsets, cached exact answers are the strongest
tested option. Affine scores are also exact on the same quantized inputs.
Their evaluated center-specific preprocessing repeats for new centers;
reusable public-center Full avoids repeated setup and is the relevant exact
changing-center control. A reusable affine parameterization remains possible
for updates restricted to a fixed span. When the complete center family is
known before preprocessing, exact family scores provide an additional
alternative (Section~\ref{sec:family-baseline}); its proof cost is unmeasured
and dominance over it is not claimed. The proposed benefit is lower total
verifiable-audit cost in supported changing-center workloads, not greater
accuracy than an exact baseline or faster individual queries than a cache.

\subsection{Representation and cost--guarantee limits}
The centroid-informed head preserves the deployed center span and fills its
remaining budget using residual PCA. It is not a trained neural encoder.
Both head and PCA receive the same reusable certificate circuit. In the new
rank 63 early-update grid their decision rates tie across all fitted
configurations, although assignment coverage differs. Historical native and
shift studies likewise do not establish uniform learned superiority.
Accordingly the contribution is representation-agnostic; application-specific
basis choice remains an externally approved approximation policy.

Increasing rank can improve coverage but raises computation and storage, and
the rounded Gram defect changes too. The new fixed-precision/nested-basis
control isolates this trade-off on one fit; older rank pilots changed
precision and batching together. The paired relation also has quadratic
center-count dependence. On the completed text proof fit it is only modestly
cheaper than Full and resolves no actual update. Paying for paired proofs and
then exact fallback can be more expensive than going directly to Full.
Network latency, service concurrency and an on-demand fallback prover remain
unmeasured.

\subsection{Fallback economics and deployment selection}
A certificate-first policy must charge failed attempts. For a fixed sequence
of requests, write $C_P$ for the paid paired pipeline and $F_j$ for the
incremental exact fallback cost of unresolved candidate $j$, including required
public verification. Let $F_0$ be a baseline proof cost shared by the exact
service and fallback, and $J$ the complete candidate set. With identical cost
boundaries, certificate-first beats direct exact evaluation only if
\begin{equation}
 C_P+F_0+\sum_{j\in U}F_j < F_0+\sum_{j\in J}F_j,
\end{equation}
or equivalently $C_P<\sum_{j\in J\setminus U}F_j$. If the paired and exact
pipelines share additional work, only measured reuse may be credited. A proof
replay time alone cannot supply the missing fallback generation cost. When no
request resolves, paying a positive paired cost before the same exact service
gives no saving under this accounting. The text example therefore favors
starting with Full; the no-fallback image deployment ratio is conditional.

A practical deployment must choose its service contract before examining
private audit outcomes: exact known-family preprocessing where appropriate,
Full for exact decisions, or certificate-first with explicit inconclusive
outcomes and a fallback budget. The present data do not establish an optimal
selection policy, an online fallback latency distribution or a production
regression-detection service. Those require targeted measurements rather than
pooling further positive early-update examples.

\subsection{Evidence and submission boundaries}
The 2,160-cell numerical grid broadens learning-problem evidence but shares
splits, centers and updates within its 15 dataset/seed replications. It is
not 2,160 independent experiments. At the stated cutoff, new paired proof
measurements cover three fitted 128-record workloads; 15 of 32 main cases are
complete, with separate portability and synthetic controls. These completed
cases are a temporal snapshot of a fixed queue, not a campaign-wide success
estimate. The main queue and additional 1,024-record paired scale queue have since
reported completion through the server console. The scale comparison is
reported separately as provisional console evidence; second-seed case files
and full scale proof exports still await recovery and independent replay. Earlier actual 1,024-record absolute
certificate proofs are retained as a distinct cohort.

The matched one-repetition deployment's 2.07-fold wall-time ratio includes
public replay, but represents one fitted workload on one resource-bounded
host. Other alternative costs reconstruct recorded components and exclude
fitting and process overhead as stated. Windows/Linux portability confirms
execution on a second environment, not a cross-host speed distribution.
Public benchmarks and algorithmic Lloyd traces do not replace a production
update study. The experiments prove neither model training nor feature
extraction, and fixed-catalog preprocessing limits predictive-learning claims.

The projection/residual identities, Cauchy--Schwarz, min/max monotonicity,
Merkle authentication and Groth16 are established. The proposed systems
contribution is their checked integer composition into reusable cost-change
certificates, independently approved tasks and a public fallback interface.
Implemented utility examples sharpen this contribution without establishing
a new general sketching theorem or unrestricted novelty-firstness.
The staged artifact is privately verifiable and has passed fresh replay, but
public accessibility remains pending. Author details, declarations and final
scientific approval remain with the authors.

\section{Conclusion}\label{sec:conclusion}
Authenticated reusable statistics support integer intervals for changes in
exact quantized clustering cost. A conclusive sign guarantees an update
decision; exact fallback resolves inconclusive tasks at additional cost.
Completed image examples and one 2.07-fold deployment saving demonstrate
conditional usefulness. The separately labelled 1,024-record console result
extends this evidence with a provisional 2.10-fold component-plus-replay
ratio and four resolved updates for both head and PCA; original export
recovery and independent replay are still required. Text, mature updates, the single-fit scale comparison and incomplete
recovery of the completed campaign limit generalization. The results do not
establish uniform learned superiority.
